\section{\bt{Background and Problem}{背景与问题定义}}
\label{sec:background}

\subsection{\bt{Shared Computations in Agent Workloads}{智能体负载中的共享计算}}
\label{sec:workload-characterization}
% English: input both tables here so Table I lands on page 2 and Table II (two-column) on page 3.
\begin{table}[t]
\caption{\bt{Characteristics of application SQL and memoryless agent workloads.}{应用查询与无记忆智能体会话的特征。}}
\label{tab:workload-characterization}
\footnotesize
\setlength{\tabcolsep}{4pt}
\renewcommand{\arraystretch}{1.08}
\begin{tabularx}{\linewidth}{@{}Lrr@{}}
\toprule
& \bt{Redset}{Redset} & \bt{Agents}{智能体}\\
\midrule
\bt{LLM calls per session}{每个会话的大模型调用} & -- & 4.6\\
\bt{\quad issuing no answer SQL}{\quad 不发出答案所用 SQL 的} & -- & 57\%\\
\bt{Session time before the answer SQL}{答案所用 SQL 之前的会话时间} & -- & 71\%\\
\addlinespace
\bt{Lookups repeating an earlier fact}{查找重复已有事实} & -- & 93\%\\
\bt{\quad share of lookup DB time}{\quad 占查找数据库时间} & -- & 77\%\\
\bt{SQL lookups repeating a fact}{SQL 查找重复事实} & -- & 38\%\\
\bt{\quad with the same SQL text}{\quad 其中 SQL 文本相同} & -- & 0\%\\
\addlinespace
\bt{SQL repeating an earlier template}{SQL 重复已有模板} & 93.4\% & 24\%\\
\bt{Templates per 100 statements}{每百条语句的模板数} & 6.6 & 76\\
\bt{Answered from the result cache}{由结果缓存回答} & 34\% & --\\
\addlinespace
\bt{SQL that fails}{出错的 SQL} & 0.35\% & 1.7\%\\
\bt{Wrong answers from SQL without error}{错误答案来自无报错 SQL} & -- & 32 / 32\\
\bt{Calls that validate}{验证类调用} & \bt{none}{无} & 4.8\%\\
\bt{Sessions checking key uniqueness}{检查键唯一性的会话} & -- & 16\%\\
\bottomrule
\end{tabularx}
\end{table}
\begin{table*}[t]
\caption{\bt{Comparison of reuse, lookup, and validity capabilities.}{复用、查找与有效性能力对比。}}
\label{tab:capabilities}
\footnotesize
\newcommand{\capyes}{\ding{51}}
\newcommand{\capno}{\ding{55}}
\setlength{\tabcolsep}{3pt}
\renewcommand{\arraystretch}{1.12}
\begin{tabularx}{\linewidth}{@{}P{0.18\linewidth}P{0.20\linewidth}*{5}{>{\centering\arraybackslash}X}@{}}
\toprule
\multicolumn{4}{c}{\textbf{\bt{Reuse}{复用}}} & \multicolumn{3}{c}{\textbf{\bt{Validity}{有效性}}}\\
\cmidrule(lr){1-4}\cmidrule(lr){5-7}
\textbf{\bh{Design}{设计}} & \textbf{\bh{Reusable object}{复用对象}} & \textbf{\bh{Knowledge\\extraction}{知识提取}} & \textbf{\bh{Ambiguity-\\aware lookup}{歧义感知查找}} & \textbf{\bh{Derived\\data checks}{派生数据检查}} & \textbf{\bh{Shared\\check results}{共享检查结果}} & \textbf{\bh{Data-update\\handling}{数据更新处理}}\\
\midrule
AgentSM~\cite{biswal2026agentsm} & \bhl{Trajectories; tools}{轨迹；工具} & \capyes & \capno & \capno & \capno & \capno\\
G-Memory~\cite{zhang2025gmemory} & \bhl{Trajectories; insights}{轨迹；洞见} & \capyes & \capno & \capno & \capno & \capno\\
DataLab~\cite{weng2025datalab} & \bhl{Business logic}{业务逻辑} & \capyes & \capno & \capno & \capno & \capno\\
TeCoD~\cite{jivani2025tecod} & \bhl{Question--SQL templates}{问题--SQL 模板} & \capyes & \capno & \capno & \capno & \capno\\
dbt Semantic Layer (MetricFlow)~\cite{dbt2026metricflow,dbt2026slcache} & \bhl{Metrics; cached results}{指标；缓存结果} & \capno & \capyes$^{c}$ & \capno & \capno & \capyes$^{a}$\\
Databricks metric views~\cite{databricks2026metricviews,databricks2026materialization} & \bhl{Metrics; materialized data}{指标；物化数据} & \capno & \capyes$^{c}$ & \capno & \capno & \capyes$^{b}$\\
vCache~\cite{schroeder2026vcache} & \bhl{LLM responses}{LLM 回答} & \capno & \capyes$^{d}$ & \capno & \capno & \capno\\
\bhl{Semantic cache~\cite{dar1996semantic}}{语义缓存} & \bhl{Data regions}{数据区域} & \capno & \capyes$^{e}$ & \capno & \capno & \capno\\
TxCache~\cite{ports2010txcache} & \bhl{Computation results}{计算结果} & \capno & \capyes$^{e}$ & \capno & \capno & \capyes\\
\midrule
\textbf{\system} & \textbf{\bhl{Metric definitions}{指标定义}} & \capyes & \capyes & \capyes & \capyes & \capyes\\
\bottomrule
\end{tabularx}
\par\smallskip
\bt{\capyes: supported; \capno: not supported in the cited design. Ambiguity-aware lookup: a natural-language request is matched without SQL, competing definitions are returned when it leaves a choice open, and absent metrics are rejected (TeCoD's NLI matcher rejects, but returns one template). Derived data checks are generated from each definition's grain, joins, and filters; dbt-style layers support uniqueness, relationship, and custom SQL tests, declared by hand.}{\capyes：支持；\capno：所引设计不支持。歧义感知查找：不合成 SQL 即可匹配自然语言请求，请求未说明某项口径时返回全部竞争定义，并拒识目录中没有的指标（TeCoD 的 NLI 匹配器可以拒识，但只返回一个模板）。派生数据检查由每个定义的粒度、连接与过滤自动生成；dbt 类语义层支持唯一性、关系与自定义 SQL 测试，但需手工声明。}
\bt{$^{a}$With declarative caching. $^{b}$With unaggregated materialization and scheduled refresh. $^{c}$Metrics are requested by their declared name. $^{d}$Statistical error bound on response text. $^{e}$Matches given queries or calls, not natural-language requests.}{$^{a}$启用声明式缓存。$^{b}$启用非聚合物化与定时刷新。$^{c}$按声明的名称请求指标。$^{d}$对回答文本的统计错误率上界。$^{e}$匹配给定的查询或调用，而非自然语言请求。}
\end{table*}


\bi{Traditional database applications put domain logic in fixed code paths that issue parameterized SQL templates~\cite{ma2018querybot}. To see how data agents differ, we compare 441 million read queries across 400 Redshift instances over three months from Redset~\cite{vanrenen2024redset} (grouped by structural fingerprint per instance) with 100 memoryless \llmmodel{} sessions answering 25 analytical questions over five metrics (521 database calls, including 176 SQL statements). Calling each schema inspection, data probe, or validation query a \emph{lookup} and grouping lookups with the same intent and target tables/columns into an \emph{inspection target}, Table~\ref{tab:workload-characterization} shows a sharp contrast. In Redset, 93.4\% of reads repeat an earlier fingerprint (6.6 templates per 100 statements), so a result cache answers 34\% of reads; in agent workloads, 93\% of lookups revisit an earlier inspection target, yet agents write 76 distinct SQL templates per 100 statements with 0\% exact SQL-text reuse across revisiting probes.}
{传统数据库应用把业务逻辑写在固定代码路径中，发出参数化 SQL 模板~\cite{ma2018querybot}。为看清数据智能体的差异，我们将 Redset~\cite{vanrenen2024redset} 中 400 个 Redshift 实例三个月内的 4.41 亿条读查询（按实例内结构指纹分组）与围绕五个指标回答 25 道分析题的 100 个无记忆 \llmmodel{} 会话（共 521 次数据库调用，含 176 条 SQL）进行对比。把每次模式查看、数据探查或验证查询称为一次\emph{查找}，把意图和目标表列相同的查找归为同一\emph{检查目标}，表~\ref{tab:workload-characterization} 呈现出鲜明对比：Redset 中 93.4\% 的读查询重复既有指纹（每百条仅 6.6 个模板），结果缓存能直接回答 34\% 的读请求；而在智能体负载中，尽管 93\% 的查找重复访问已有检查目标，智能体每百条语句却写出 76 个不同 SQL 模板，重复探查间的精确 SQL 文本复用率为 0\%。}

\bi{This difference arises because each agent session mixes a \emph{shared domain computation}---fact grain, join path, business filter, and measure formula---with \emph{request-specific parameters} (reporting periods, rankings, or period comparisons) and \emph{implicit data assumptions}. Rediscovering that shared computation dominates both cost and mistakes: 57\% of LLM calls issue no final-answer SQL, 71\% of session wall-clock time passes before the answer SQL runs (consistent with Liu et al., who report 12.67 SQL statements per task on o3 and a 27.7\% drop in column exploration from expert hints~\cite{liu2026agents}), and all 32 wrong answers come from SQL that executes without error---31 on questions that name a metric without stating its formula. Meanwhile, only 16\% of sessions check grain-key uniqueness. Predictive modeling agents show the same pattern, repeatedly probing entity joins, row grains, and look-back windows (such as 90-day customer spend) across tasks and time cutoffs (Sections~\ref{sec:features} and~\ref{sec:relbench-eval}).}
{出现这一差异的原因在于，每次智能体会话都把\emph{共享的领域计算}（事实粒度、连接路径、业务过滤与度量公式）同\emph{单次请求的参数}（报告期间、排序或跨期对比）以及\emph{隐含的数据假设}混在一起。重复摸索共享计算主导了开销与错误：57\% 的大模型调用未发出最终答案 SQL，71\% 的会话耗时发生在该 SQL 运行之前（与 Liu 等人的观察一致：o3 每题平均执行 12.67 条 SQL，提供专家提示可减少 27.7\% 的列探查~\cite{liu2026agents}）；全部 32 个错误答案都来自未报错的合法 SQL，其中 31 个出现在只给指标名称而未给公式的问题上；同时仅 16\% 的会话检查了粒度键唯一性。预测建模智能体也呈现同样规律，在不同任务与时间截点之间反复探查相同的实体连接、聚合粒度与回看窗口（如客户近 90 天消费额；见第~\ref{sec:features}、\ref{sec:relbench-eval} 节）。}

\subsection{\bt{Design Space of Existing Systems}{现有系统的设计空间}}
\label{sec:alternatives}

\bi{Table~\ref{tab:capabilities} groups existing systems by \emph{Reuse} (stored object, session extraction, and ambiguity-aware lookup) and \emph{Validity} (derived data checks, shared outcomes, and update maintenance). Agent memory frameworks store past trajectories per agent~\cite{packer2023memgpt,chhikara2025mem0,xu2025amem,zhao2024expel,wang2025awm,zhang2025apc,zhang2026ace} or across agents~\cite{zhang2025gmemory,fioresi2026learning,biswal2026agentsm,databricks2026memcurrent}; even with invalidation contracts~\cite{wu2026invalidation}, transactional belief checks~\cite{li2026memtx}, or provenance tracking~\cite{wang2026pacmi}, notes and workflows summarized by prompts compress text without checking SQL structure or replaying answers, so filters or joins can be lost silently. Conversely, script and example stores such as DataLab~\cite{weng2025datalab}, Snowflake VQR~\cite{snowflake2026vqr}, and TeCoD~\cite{jivani2025tecod} keep concrete scripts or question--SQL pairs that pile up across slices and still need LLM editing at reuse, whereas semantic layers (MetricFlow~\cite{dbt2026metricflow,dbt2026metricflowvalidation,dbt2026slcache}, Databricks metric views~\cite{databricks2026metricviews,databricks2026materialization}) and query caches~\cite{dar1996semantic,ports2010txcache} serve hand-written metrics or cached query regions rather than extracting definitions from agent sessions.}
{表~\ref{tab:capabilities} 按\emph{复用}（存储对象、会话提取与歧义感知查找）和\emph{有效性}（派生检查、结果共享与更新维护）归纳现有系统。智能体记忆框架在单智能体内~\cite{packer2023memgpt,chhikara2025mem0,xu2025amem,zhao2024expel,wang2025awm,zhang2025apc,zhang2026ace}或跨智能体~\cite{zhang2025gmemory,fioresi2026learning,biswal2026agentsm,databricks2026memcurrent}保存历史轨迹；即便引入失效契约~\cite{wu2026invalidation}、事务信念校验~\cite{li2026memtx}或溯源追踪~\cite{wang2026pacmi}，靠提示词总结出的文本笔记与工作流并未检查 SQL 结构或回放核对答案，仍可能悄悄丢失过滤或连接。另一边，DataLab~\cite{weng2025datalab}、Snowflake VQR~\cite{snowflake2026vqr} 与 TeCoD~\cite{jivani2025tecod} 保存具体脚本或问题--SQL 示例，随切片不断累积且复用时仍需大模型改写；企业语义层（MetricFlow~\cite{dbt2026metricflow,dbt2026metricflowvalidation,dbt2026slcache}、Databricks metric views~\cite{databricks2026metricviews,databricks2026materialization}）与查询缓存~\cite{dar1996semantic,ports2010txcache} 则服务于人工编写的指标或缓存结果，不能从智能体会话自动提取定义。}

\bi{Because neither example banks nor prompt summaries expose a structured relational definition, lookup and maintenance suffer as well. Embedding caches~\cite{bang2023gptcache}, threshold-calibrated vCache~\cite{schroeder2026vcache}, production semantic caches~\cite{rathore2026cogni,singh2026semanticcache}, and entailment-matched TeCoD~\cite{jivani2025tecod} return at most one entry per request---missing competing definitions on ambiguous queries and risking false hits on absent metrics unless an LLM scans the catalog online. Likewise, TTLs and schema hashes~\cite{singh2026semanticcache} miss data-only breaks under fixed schemas, while data validation suites~\cite{song2021autovalidate,breck2019validation,chen2026prismadv} check single-table batches rather than maintaining multi-table definitions across concurrent agents. \system\ connects all three stages through checked executable definitions.}
{由于示例库和提示词摘要都没有给出结构化的关系定义，后续查找与维护也随之受限：向量缓存~\cite{bang2023gptcache}、vCache~\cite{schroeder2026vcache}、生产级语义缓存~\cite{rathore2026cogni,singh2026semanticcache}与 TeCoD~\cite{jivani2025tecod} 每次最多返回一条结果，既无法在含糊请求上返回多条竞争定义，若不在线调用大模型扫描全目录又容易在缺失指标上误命中；过期时间与模式哈希~\cite{singh2026semanticcache}发现不了表结构不变的数据违例，单表数据校验工具~\cite{song2021autovalidate,breck2019validation,chen2026prismadv}也无法跨并发智能体维护多表定义。\system{} 通过经核验的可执行定义把三阶段串联起来。}

\subsection{\bt{Problem Definition}{问题定义}}
\label{sec:problem-formulation}

\bi{We model cross-agent knowledge sharing over a changing database as constrained cost minimization. Starting from an empty catalog, $N$ agents answer $n$ requests interleaved with database writes. Request $j$ arrives as natural-language text $\rho_j$ asking for one of $k$ underlying computations (an analytical metric or a predictive feature set) with runtime parameters $p_j$ (a reporting period, slice, or seed entities and cutoffs) on database snapshot $s_j$, whose unobserved ground-truth output is $a^*_j$. To answer $\rho_j$, an agent either explores the database from scratch (Steps~1--4 in Figure~\ref{fig:metric-record}) at mean cost $c_e$, or reuses a published definition by compiling it with $p_j$ on $s_j$ to produce $\hat a_j$ at mean runtime cost $c_r$ (covering lookup and snapshot checks). Let $c_\ell$ denote the cost of extracting, checking, and publishing a definition, and $M(\pi)$ the total maintenance cost of rechecks and repairs across writes. A sharing policy $\pi\in\Pi$ chooses the published set $\mathcal P$, each request's reuse decision $y_j\in\{0,1\}$, and the maintenance actions.}
{我们将变化数据库上的跨智能体知识共享建模为带约束的成本最小化问题。从空目录出发，$N$ 个智能体依次回答穿插数据库写入的 $n$ 个请求。请求 $j$ 以自然语言文本 $\rho_j$ 给出，指向 $k$ 类底层计算（分析指标或预测特征集）之一及运行期参数 $p_j$（报告期间、切片，或种子实体与时间截点），读取数据库快照 $s_j$，其运行期不可观测的真实输出记为 $a^*_j$。回答 $\rho_j$ 时，智能体要么以平均开销 $c_e$ 从头探查数据库（图~\ref{fig:metric-record} 步骤 1--4），要么以平均运行期开销 $c_r$（含查找与快照检查）复用已发布定义，在 $s_j$ 上代入 $p_j$ 编译执行得到 $\hat a_j$。令 $c_\ell$ 表示提取、核验并发布一份定义的开销，$M(\pi)$ 表示数据写入引发的重检与修复总开销。共享策略 $\pi\in\Pi$ 决定发布集合 $\mathcal P$、每个请求的复用决策 $y_j\in\{0,1\}$ 以及维护动作。}

\begin{problem}[\bt{Shared reuse under data updates}{数据更新下的共享复用}]
\label{prob:shared-reuse}
\bid{Given a request stream, interleaved writes, and an error tolerance $\alpha$, find a sharing policy $\pi\in\Pi$ that solves}
{给定请求流、穿插的数据写入与错误容忍度 $\alpha$，求解共享策略 $\pi\in\Pi$：}
\begin{align}
\min_{\pi\in\Pi}\quad & \sum_{j=1}^{n}\bigl[(1-y_j)\,c_e+y_j\,c_r\bigr]+c_\ell\,\lvert\mathcal P\rvert+M(\pi) \label{eq:objective}\\
\textnormal{s.t.}\quad & \sum_{j=1}^{n} y_j\,\mathbb{1}\bigl[\hat a_j\neq a^*_j\bigr]\;\le\;\alpha\sum_{j=1}^{n} y_j. \label{eq:correctness}
\end{align}
\end{problem}

\bi{Exploring every request (${y_j=0}$, ${\mathcal P=\varnothing}$) is always valid at cost $n\,c_e$. Because definitions come only from verified exploration sessions, any policy needs at least $k$ explorations (${\sum_j(1-y_j)\ge k}$), whereas private per-agent memories need up to $Nk$. Approaching the $k$-exploration bound while keeping $\lvert\mathcal P\rvert$, $c_r$, and $M(\pi)$ small under Constraint~\eqref{eq:correctness} motivates \system{}'s three mechanisms: (1)~compressing sessions into checked definitions that merge duplicates to keep $\lvert\mathcal P\rvert$ small and compile across $p_j$ without \emph{lossy extraction} (C1, Section~\ref{sec:publication}); (2)~matching requests against catalog choices to raise safe reuse ${y_j=1}$ at low $c_r$ without \emph{false hits} (C2, Section~\ref{sec:lookup}); and (3)~skipping unchanged tables, sharing check results, repairing safely, and executing in one snapshot to keep $M(\pi)$ low without \emph{stale answers} (C3, Section~\ref{sec:maintenance}).}
{对每个请求都从头探查（$y_j=0,\mathcal P=\varnothing$）总是可行，总开销为 $n\,c_e$。由于定义只能来自完成探查并核验的会话，任何策略至少需要 $k$ 次探查（$\sum_j(1-y_j)\ge k$），而单智能体私有记忆最多需要 $Nk$ 次。要在逼近 $k$ 次探查下界的同时压低 $\lvert\mathcal P\rvert$、$c_r$ 与 $M(\pi)$ 并满足错误率约束~\eqref{eq:correctness}，自然引出 \system{} 的三项机制：（1）将会话压缩为经核验的定义并合并重复发现，控制 $\lvert\mathcal P\rvert$ 且跨参数 $p_j$ 准确编译，消除\emph{有损提取}（C1，第~\ref{sec:publication} 节）；（2）按目录口径匹配请求，以低开销 $c_r$ 提高安全复用率 $y_j=1$，避免\emph{误命中}（C2，第~\ref{sec:lookup} 节）；（3）跳过未变更表、共享检查结果、安全修复并在同一快照内执行，控制 $M(\pi)$ 并杜绝\emph{陈旧错答}（C3，第~\ref{sec:maintenance} 节）。}
